\title{基于程序化内容生成与实时图形渲染的算法设计}
\author{杨其臻}
\date{Nov 15, 2026}
\maketitle
在开放世界与元宇宙的开发实践中，场景规模从数十平方公里攀升至数千平方公里时，手工建模的成本瓶颈与重复感问题愈发突出。程序化内容生成（PCG）与实时渲染管线在这一背景下形成交汇：前者负责“有什么”，后者解决“怎么看”。本文将围绕这一交汇点，系统梳理算法选型、管线协同与性能权衡，帮助读者建立可落地的技术认知。\par
\chapter{理论基石}
\section{程序化内容生成算法综述}
在生成连续且自然的标量场时，经典噪声函数是最基础的工具。值噪声通过在格点上随机赋值并双线性插值实现平滑过渡；梯度噪声（Perlin）则在格点存储梯度向量，利用点积与缓和曲线构造平滑结果；simplex 噪声进一步将立方体格网替换为单纯形，降低维度依赖并提升各向同性。时间复杂度上，三者均为常数级，但在高维或高频调用时，缓存命中率与 SIMD 向量化成为关键瓶颈。\par
分形叠加可使噪声在不同尺度上同时呈现细节。fBm（fractional Brownian motion）通过多层噪声按指数衰减振幅叠加，得到连续且具有自相似性的高度场；ridged multifractal 则在每层取绝对值并反转，突出山脊；worley noise 以距离场为基础，适合模拟细胞或碎石纹理。空间复杂度随层数线性增长，因此在 GPU 上通常采用纹理预计算或 Compute Shader 动态生成。\par
基于规则与搜索的结构化生成则面向离散对象。L-System 通过符号重写迭代产生树枝或城市道路；Wave Function Collapse（WFC）在约束满足框架下从样本图块中“坍缩”出全局一致的布局；遗传算法则在多目标优化场景下搜索参数空间。三者均需在生成阶段进行剪枝或约束传播，以避免指数级回溯。\par
\section{实时渲染管线中的 PCG 数据流}
在实时渲染中，PCG 生成的数据必须满足资源预算与 LOD 策略。几何预算限制每帧三角形数量，纹理流送则依据视锥体与 LOD 选择性解码；DrawCall 合并通过间接绘制与 Mesh Shader 路径，将海量图元打包为少量调用，从而降低 CPU 开销。\par
GPU-Driven 管线进一步把生成逻辑迁移至 Compute Shader：首先在显存中分配高度场或实例参数缓冲区，随后启动核函数并行计算顶点与索引；最后通过 Indirect Draw Buffer 驱动 Mesh Shader，无需 CPU 读回。此路径的关键在于缓冲区对齐与内存屏障，确保生成与绘制在同一帧内完成。\par
\chapter{核心算法设计}
\section{地形与自然环境}
地形高度场生成通常采用多层噪声叠加再叠加侵蚀模拟。首先生成由 fBm 构成的初始高度场，其分辨率与世界尺寸成正比；随后在 GPU 上运行 Hydraulic Erosion 迭代：每帧计算通量、溶解与沉积，时间复杂度为图元数乘以迭代次数。Thermal Erosion 则在相邻像素间按坡度阈值做质量迁移，模拟风化崩塌。两类侵蚀均需在 Compute Shader 中利用共享内存缓存邻域，减少全局内存带宽。\par
植被分布采用 Poisson-Disk Sampling 保证最小间距，再叠加生物群落规则。采样过程可通过 Compute Shader 中的 Grid Jitter 实现线性时间复杂度；群落规则则读取土壤湿度与坡度纹理，输出实例变换矩阵。最终实例数据以 Mesh Shader 的 per-meshlet 形式提交，避免 CPU 侧的巨大缓冲区拷贝。\par
实时 LOD 采用 CDLOD（Continuous Distance-Dependent LOD）或 Geoclipmap 方案。CDLOD 将地形划分为嵌套规则网格，每层对应一个 mipmap 级别；GPU 根据视点距离选择顶点层级，顶点着色器仅需一次采样即可完成位移。Geoclipmap 则在环形带上维护多分辨率高度图，通过异步 Compute Shader 增量更新，带宽需求与环宽成正比。\par
\section{建筑与人工结构}
基于语法的建筑立面生成可借助 CGA Shape Grammar。在初始符号“Lot”上依次应用 Split、Repeat、Component Split 规则，生成墙面、窗户与阳台层级。每条规则可附加概率或参数，使立面在保持结构规律的同时呈现多样性。求值过程在 CPU 端进行符号展开，再将最终图元批量上传 GPU；若需动态修改，可将符号栈编码为 Compute Shader 的任务队列。\par
结构稳定性约束通过物理启发式剪枝实现。对每条新生成的梁或柱，立即执行简化静力分析：若其应力超出阈值，则删除或加粗截面。网格加固则在拓扑阶段插入对角撑或剪力墙，降低后续物理模拟的迭代次数。该过程可视为在符号展开树上的一次剪枝遍历，时间复杂度与符号数量呈线性关系。\par
\section{材质与纹理}
过程化材质网络沿用 Substance 的节点范式，但在运行时映射为自定义 Shader Graph。每个节点对应一个或多个 GLSL/HLSL 函数，输入输出均为四元向量或标量；常量折叠与公共子表达式消除可在编译期完成，从而降低运行时指令数。Graph 的有向无环结构天然支持多线程求值，适合在 Async Compute Queue 中与主渲染并行。\par
纹理合成与运行时流送依赖虚拟纹理（VT）与 Tile 按需解码。VT 将巨型纹理划分为固定大小的物理页，着色器通过页表间接寻址；当发生缺页时，CPU 或后台 Compute Shader 负责解码对应 Tile 并上传。解码可选用 ASTC 或 BC7 自适应压缩，在保证画质的前提下将带宽降至原图的 10\%{} 以下。初始加载仅需索引表与少量热页，因此可在 200 ms 内完成 10 万栋建筑的材质准备。\par
\chapter{性能与工程实践}
\section{内存与带宽}
几何数据压缩在顶点层面采用量化：位置坐标从 float32 降至 int16，法线与切线编码为 octahedral 映射后的 int8，整体压缩率可达 4:1。Meshlet 结构进一步将 128 个顶点与 256 个三角形打包为缓存友好的块，Mesh Shader 可直接在 SM 内完成剔除与光栅化，减少 DrawCall 数量。空间复杂度从 O(n) 降至 O(n/128)，同时提升缓存命中率。\par
纹理与贴图预算则依赖 ASTC/BC7 自适应压缩与 UDIM 虚拟纹理。ASTC 支持可变块大小，可在 4×4 到 12×12 之间权衡质量与带宽；BC7 则在 DirectX 平台提供固定 4×4 块的高质量压缩。UDIM 通过多 UV 岛将纹理拆分到多张 8K 页，仅在视锥体内解码对应页，从而把峰值显存控制在 2 GB 以内。\par
\section{多线程与异步管线}
Job System 与帧图（Frame Graph）负责解析 PCG Job 与 Render Job 的依赖。帧图以资源为节点、Pass 为边，自动插入屏障与别名分配；PCG Job 在 CPU 端生成任务描述符，再由 Job System 调度到 Async Compute Queue。依赖解析的时间复杂度与 Pass 数量呈线性，实际开销可通过拓扑排序的缓存友好实现控制在 1 ms 以内。\par
增量更新与脏矩形机制用于降低噪声重采样成本。仅当视点移动超过阈值或参数变化时，才在脏矩形区域内重新生成噪声；增量 GPU Upload 则利用 ring buffer 与 fence，将更新数据与绘制解耦。时间复杂度从全图 O(n) 降至 O(d)，其中 d 为脏矩形面积，通常占全图的 5\%{} 以下。\par
\section{质量-性能权衡}
偏差-方差控制需在噪声带宽与 TAA 之间取得平衡。过低带宽导致高频噪声被 TAA 拖影，过高则产生摩尔纹；解决方案是在噪声生成阶段加入可控的哈希抖动，并在 TAA 阶段根据局部方差自适应调整混合系数。实验表明，在 4K 分辨率下，抖动幅度取 0.5 pixel 时，收敛速度与画质最佳。\par
运行时预算监控可借鉴 Nanite 的 LOD 策略：每帧统计三角形与纹理带宽，若超出阈值则在下一帧动态降低噪声层数或提高压缩率。自适应质量标尺以毫秒为单位度量 PCG Job 与 Render Job 的耗时，利用 PID 控制器调整质量参数，使帧时间稳定在 16.6 ms 附近。\par
\chapter{案例研究}
\section{开放世界地形 demo}
该 demo 覆盖 256 km² 地形，在 4K 分辨率下维持 30 fps。全流程在 GPU 上完成：Compute Shader 生成 8K×8K 的高度场与法线图，耗时 3.2 ms；Mesh Shader 驱动 CDLOD 渲染，平均每帧 2.1 百万三角形，剔除与光栅化耗时 8.7 ms；Hydraulic Erosion 每 10 帧更新一次，单次 4.8 ms。显存峰值 1.8 GB，主要来自虚拟纹理与 Meshlet 缓冲区。\par
\section{程序化城市}
程序化城市包含 10 万栋建筑，初始加载时间小于 200 ms。CGA Shape Grammar 在 CPU 端并行展开符号树，输出 1200 万个图元；Mesh Shader 路径将图元打包为 9 万个 meshlet，DrawCall 数量降至 120 次。纹理流送使用 4K 虚拟纹理与 ASTC 压缩，峰值带宽 1.2 GB/s；Job System 将 PCG 与渲染解耦，使主线程帧时间稳定在 11 ms。\par
\chapter{工具链与调试}
可视化编辑器采用节点式噪声编辑器，节点可实时在 GPU 上求值并预览。热重载通过内存映射文件实现：编辑器将节点参数写入共享内存，运行时守护线程检测到变化后立即重建 Shader 变体，无需重启引擎。预览窗口支持 3D 漫游与参数动画，便于快速定位噪声异常。\par
性能分析依赖 RenderDoc 与 GPU Trace。RenderDoc 可捕获 Compute Shader 的中间缓冲区，检查噪声生成是否出现 NaN 或带宽溢出；GPU Trace 则以 ns 级精度统计每条指令的占用率，定位到 fBm 中最内层循环的 40\%{} 占用瓶颈。结合脏矩形可视化，可直观看到增量更新是否覆盖全部必要区域。\par
\chapter{未来展望}
AI 协同正成为 PCG 的新方向。生成对抗网络（GAN）可学习真实地形与立面的分布，从而自动搜索噪声参数或文法规则；强化学习则在运行时根据玩家行为动态调整生成策略，实现“常玩常新”。硬件演进方面，Geometry Pipeline 2.0 与下一代 Mesh Shader 将进一步降低 CPU 负担，把更多生成逻辑迁移至固定功能单元。\par
结语：程序化内容生成与实时渲染的深度融合，已从实验室走向工业实践。算法选型需兼顾数学特性与硬件约束，管线设计需以异步与增量为核心，性能权衡则依赖实时监控与自适应策略。未来，随着 AI 与新硬件的加入，这一领域仍将持续演进，为开放世界与元宇宙提供更丰富、更高效的内容供给。\par
